\section{Coherent and second-order absorption}\label{app:E.26}

An axion-like field of mass \SI{1}{\micro eV} stays coherent for $t_{\mathrm{coh}} = 1/(m_\phi v_0^2) \approx \SI{1.04}{ms}$, with $v_0$ the circular speed of the standard halo (\cref{sec:1.2}), and over $1/(m_\phi v_0) \approx \SI{249}{m}$, so a nuclear-spin sample or an acoustic resonator sees one classical drive, not a stream of quanta. A fermion of \SI{100}{keV} absorbed in germanium deposits nearly its whole mass at a momentum transfer of the same size. A dark photon can be absorbed in pairs, or scattered into a photon. These processes lie outside the single-vertex golden rule of \crefrange{app:E.14}{app:E.16}, and \cref{sec:10.8,sec:10.9,sec:11.8,sec:14.6} derive them; what they need from the material is mostly a response that the methods of \cref{app:E.4,app:E.6,app:E.7} already compute, read on a different line of the $(q, \omega)$ plane.

\subsection{Coherent-regime material inputs}

The first decision is the regime, and it is made before any material calculation. The target responds on a time $1/\varGamma_T$, which is $T_2$ for spins and the mode quality over the mode frequency for a resonator. For an integration time $t_{\mathrm{int}}$ long compared with the smaller of $t_{\mathrm{coh}}$ and $1/\varGamma_T$, the mean absorbed power equals the velocity-averaged golden rule, and the rate tables of \crefrange{chap:30}{chap:32} apply. A single record matches that mean only for $t_{\mathrm{int}} \gg t_{\mathrm{coh}}$. Below $t_{\mathrm{coh}}$ the drive is a single Rayleigh-distributed amplitude, the signal is the windowed linear response of \cref{sec:14.6}, and the sensitivity scales as $t_{\mathrm{int}}^{-1/2}$ instead of $(t_{\mathrm{int}}t_{\mathrm{coh}})^{-1/4}$. The material inputs differ by platform.

For nuclear-spin precession, the chemical shielding and electric-field-gradient tensors of every nuclear site give the inhomogeneous line as a site sum, and come from the gauge-including projector augmented-wave method~\cite{pickard2001all}. The effective electric field at the nucleus that converts a nucleon electric dipole moment into an energy shift comes from density-functional theory, and $T_1$ and $T_2$ come from data. Prepolarized samples are out of equilibrium, with $T_1$ of hours, so their spin noise is evaluated at the spin temperature fixed by the polarization $p$ through $\tanh(\beta_s\omega_{\mathrm{Lar}}/2) = p$, with $\omega_{\mathrm{Lar}}$ the Larmor frequency, not at the lattice temperature~\cite{budker2014proposal}. Bulk acoustic resonators need the elastic and piezoelectric tensors from density-functional perturbation theory (\cref{app:E.4}) and the mode frequencies and shapes from a finite-element solution of the body. They also need the orientation factor of the field in the sample (1 normal to a slab, $3\varepsilon_0/(\varepsilon_0 + 2)$ for a sphere, $\varepsilon_0$ parallel to a slab; \cref{sec:11.8}), and the measured mode quality~\cite{trickle2025piezoelectric}. For resonant-mass detectors driven by a scalar that changes $\alpha$ or $m_e$, the inputs are $\dd\ln a/\dd\ln\alpha$ and $\dd\ln a/\dd\ln m_e$ of the lattice constant $a$~\cite{arvanitaki2016sound}. The Coulomb Hamiltonian with infinite nuclear masses depends on $\alpha$ and $m_e$ only through the length unit $1/(\alpha m_e)$, so a nonrelativistic density-functional calculation with scaled constants must return $-1$ for both to numerical precision; that is the benchmark. A fully relativistic calculation of a new crystal is accepted once it reproduces the relativistic and finite-nuclear-mass corrections of copper and silicon given in \cref{sec:11.8}~\cite[Table~II]{pasteka2019material}. For magnetoelectric multiferroics, linear spin-wave theory or time-dependent density-functional theory with spin--orbit coupling on a magnetic structure computed at $T = 0$ gives the response only far below the ordering temperature. Near it, where the proposed signal peaks, a finite-temperature spin-dynamics method is needed, which \crefrange{chap:23}{chap:29} do not supply.

The frequency grid resolves the one-sided dark-matter line, $\omega \in [m_\phi, (1 + \num{3.5e-6})m_\phi]$, with at least 3000 points, and $t_{\mathrm{int}}$ runs from $10^{-2}$ to $10^3\,t_{\mathrm{coh}}$. The Asimov test statistic of \cref{sec:14.6} must change by less than \SI{1}{\percent} when the frequency step is halved, and the resonator mode frequencies must be converged to \SI{0.1}{\percent}. Two pitfalls are statistical rather than material: the random amplitude is either in the likelihood or in a penalty factor, never in both~\cite{centers2021stochastic}, and an analysis with fixed amplitude overestimates the reach, by a factor 1.99 for the median limit of a record shorter than $t_{\mathrm{coh}}$ and 2.7 for data at the detection threshold (up to 10 with other priors; \cref{sec:14.6}). Once the responses are tabulated, a point in $(m_\phi, t_{\mathrm{int}})$ takes seconds. The gauge-including projector-augmented-wave (GIPAW) and DFPT inputs use the shared setting of \cref{app:E.17}, converged to \SI{1}{\percent} in cutoff and $\vect{k}$ grid, and cost $10^{2}$ to $10^{4}$ core-hours per target, the multiferroic response about $10^{3}$.

\subsection{Fermionic absorption}

A fermion $\chi$ absorbed by the target turns into a light neutral fermion $\tilde\nu$ (\cref{sec:10.9}), and for $v \to 0$ and a massless $\tilde\nu$ the target is probed along the line $\omega = m_\chi - q$, at $q \approx m_\chi$. For $m_\chi$ from $\Eg$ to \SI{1}{MeV} this reaches momentum transfers far above the valence scale, into the cores. We need $\dd\mathcal{R}/\dd\omega$ for vector, scalar and axial operators. The scalar operator needs only the density row $S(q, \omega)$. The vector operator needs the density row and the transverse current row of $S_{ab}$: continuity turns the longitudinal current into $(\omega/q)$ times the density, so that the density enters with weight $m_\chi^2/q^2$ and the transverse current with weight one. A light mediator multiplies the scalar rate by $(m_{\mathrm{med}}^2 - 2m_\chi\omega + m_\chi^2)^{-2}$, normalized to its contact value, where $m_{\mathrm{med}}$ is the mediator mass, since the four-momentum transfer is set by the relativistic $\tilde\nu$ leg and not by $-q^2$; its expansion, $1 - 2(m_\chi^2 - 2m_\chi\omega)/m_{\mathrm{med}}^2 + \dots$, is the correction of relative order $(m_\chi/m_{\mathrm{med}})^2$ attached to the contact limit in Ref.~\cite{cox2024absorption}, with an $\omega$ dependence that tilts the spectrum.

The responses are all-electron. In crystals they come from the plane-wave calculation of \cref{app:E.7} with augmentation reconstructed to high $q$ (\cref{app:E.2}); for $q \gtrsim m_e/3$, where the core shells dominate, from Dirac--Hartree--Fock orbitals of the atom. The grid is about 200 points in $q$ per mass and 40 masses per decade. The electron spectra must change by less than \SI{2}{\percent} when the augmentation is enlarged, and the nuclear line, an impulse response broadened by the zero-point momentum $\langle p_\kappa^2\rangle$ from density-functional perturbation theory, must have its width converged to \SI{1}{\percent}. The spin convention is the commonest error: the atomic ionization form factor that the literature uses counts both spins in its sum over occupied states~\cite[Eqs.~(A.18)--(A.21)]{essig2016direct}, and a table that does not state its convention in the header can be off by a factor of two. The calculation also sizes a defect of the literature formula for the vector operator, which multiplies the ionization form factor by the free-electron matrix element instead of using continuity~\cite[Eq.~(9)]{dror2021absorption}. At bound kinematics the two differ by an $O(1)$ term $-\omega^2/m_\chi^2$. \Cref{sec:10.9} bounds the defect from one side: the current-conserving rate exceeds the free-vertex form by at least the factor $1/r$, 1.7 for the xenon $1s$ shell at $m_\chi = \SI{56.5}{keV}$, the transverse current can only raise it further, and only the relativistic corrections to the bound states, of relative size $(Z\alpha)^2 = (54/137.04)^2 = 0.155$ for the xenon K shell, have an unknown sign. Computing the transverse current and the relativistic corrections with Dirac--Hartree--Fock states fixes the size of the defect. The cost is $10^{4}$ to $10^{5}$ core-hours per crystal; the nuclear part is negligible.

\subsection{Pair, Compton and Barkas responses}

Pair absorption, dark Compton scattering and the second-order Barkas term need responses at two insertions of the field (\cref{sec:10.8}). Contact pair absorption reduces to single absorption at $\omega \approx 2m_\phi$ and $k \approx 0$, so it uses the absorption tables of \cref{app:E.16} (charge-like rows) or, for the pseudoscalar coupling, the axioelectric (spin-current) row of \cref{eq:10.7-alp} at $\omega = 2m_\phi$ (\cref{eq:10.8-P}), contracted with the pair tensor of \cref{sec:10.8}~\cite[Eqs.~(5)--(8)]{bauer2026dark}. The pseudoscalar route inherits the caveats of that row: the measured $\sigma_1$ belongs to singlet pairs while the row creates triplets, spin--orbit coupling spoils the factor $\frac34$ (\cref{sec:10.8}), and in a superconductor the spin-current conductivity replaces $\sigma_1$ only under the conditions of \cref{sec:15.4-spin}. The tests below accept the pair tensor and its velocity moments, not the row. For the scalar and pseudoscalar couplings the tensor is constant; for the vector and axial couplings it carries the relative momentum of the two absorbed particles, and its velocity moments are computed over the product of the two velocity distributions on a $20^3$ quadrature. Two tests accept it: for the scalar coupling the pair rate reduces to the single-absorption rate at $\omega = 2m_\phi$, $k \to 0$, times the constant tensor, and the velocity moments reproduce their analytic values for the standard halo to $10^{-3}$. Bauer et al.\ take a local density of \SI{0.39}{GeV.cm^{-3}}~\cite{bauer2026dark}; the rate scales as $\rho_\chi^2$, so its curves move down by $(0.3/0.39)^2 = 0.59$ for the \SI{0.3}{GeV.cm^{-3}} of \cref{sec:1.2}. Dark Compton scattering of a vector of mass $m_\phi$ on electrons bound with energies small compared with $m_\phi$ needs the shell Compton profiles, or the high-$q$ density response. It also needs the refractive index $n_r(\omega')$ of the medium at the outgoing photon energy $\omega' \approx m_\phi$; the grid is 100 points in $\omega' \in m_\phi(1 \pm 0.1)$ and 32 angles. For a free electron at rest the rate must reproduce the Thomson-limit form $\nd_e\zeta^2\sigma_{\mathrm{Th}}\,\omega'/[m_\phi(1 + \omega'/m_e)]$, with $\sigma_{\mathrm{Th}} = e^4/(6\pi m_e^2)$ the Thomson cross-section and $\omega' + \omega'^2/2m_e = m_\phi$, which agrees with the tree-level result~\cite[Eq.~(2)]{hochberg2022impact} through second order in $m_\phi/m_e$; this is the acceptance test. The Barkas term needs the quadratic density response of the electron gas, tabulated for 50 densities with $r_s$ from 1 to 6 and mapped onto the crystal by the local-density approximation on the $(q,\omega)$ grid of the loss functions of \cref{app:E.16}. It is accepted when the electron gas returns the known $Z_1^3$ (Barkas) correction to the stopping of slow charges and the linear response when the quadratic kernel is switched off. The Barkas table is new code of about $10^2$ core-hours; the other two calculations take minutes.
